% Source: https://artofproblemsolving.com/community/c3h1653434_confused_about_a_certain_symbol
% Stable identity: aops-topic-1653434-q01
\textbf{(a)} The choices are $AB,AC,AD,AE,BC,BD,BE,CD,CE,DE$, so $\dbinom52=10$.

\textbf{(b)} Count the ordered selections of $k$ objects from $n$ in two ways. Directly there are $n(n-1)\cdots(n-k+1)=\dfrac{n!}{(n-k)!}$ of them. Alternatively, first choose the set of $k$ objects ($\dbinom nk$ ways) and then order it ($k!$ ways). Hence $\dbinom nk\cdot k!=\dfrac{n!}{(n-k)!}$, which is the formula. Then
\[
\dbinom{10}{4}=\dfrac{10\cdot9\cdot8\cdot7}{4\cdot3\cdot2\cdot1}=\boxed{210}.
\]
(In (a), $\dfrac{5\cdot4}{2}=10$ agrees with the list.)
